% Annotation in Slovak
\thispagestyle{empty}

\vspace*{\fill}

\section*{Anotácia}

\begin{minipage}[t]{1\columnwidth}
	\FIITuniversitySK

	\FIITfacultySK

	Študijný program: \FIITstudyProgramSK\\

	Autor: \FIITauthor

	\FIITthesisSK: \FIITtitleSK

	Vedúci diplomovej práce: \FIITsupervisor

	\FIITdateSK
\end{minipage}

\bigskip{}

Táto práca skúma návrh a implementáciu riešení využívajúcich
dôkazy s nulovou znalosťou (ZKP) v rámci ekosystému Ethereum.
Hlavným cieľom bolo za pomoci ZkVM odbremeniť nódy na exekučnej
vrstve o potrebu znova-vykonania nového bloku. Na docielenie
bola použitiá lattice-based kryptografia, ktorá na rozdiel od eliptických
kriviek, je post-kvantovo bezpečná a na rozdiel od haší, produkuje dôkazy
s veľkosťou v desiatkach kilobajtoch a nie stovkach až tisícoch.

V práci sa podarilo navrhnúť architektúru a čiastočne implementovať
prototyp lattice-based ZkVM. Implementovaný prototyp na porovnaniach výpočtu 100. Fibonacciho čísla
vygeneroval dôkaz za 688,24 za využitia 12,34 GB RAM. V dlhšom EVM experimente
systém bežal približne 9 hodín a využil 13,7 GB RAM. Výsledky ukazujú, že
prístup je zatiaľ pomalší než etablované ZkVM riešenia, no súčasne naznačujú,
že negatívny výsledok nie je spôsobený principiálym problémom, ale
potrebuje optimalizacie. Čo sa týka veľkosti post-kvantovo bezpečných dôkazov,
tak tie sú 4 až 5krát menšie, ako porovnané riešenia.

\newpage{}\thispagestyle{empty}\medskip{}

% Annotation in English
\thispagestyle{empty}

\vspace*{\fill}

\section*{Annotation}

\begin{minipage}[t]{1\columnwidth}
	\FIITuniversity

	\FIITfaculty

	Degree Course: \FIITstudyProgram\\

	Author: \FIITauthor

	\FIITthesis: \FIITtitle

	Supervisor: \FIITsupervisor

	\FIITdate
\end{minipage}

\bigskip{}

This thesis examines the design and implementation of solutions using
Zero-Knowledge Proofs (ZKPs) in the Ethereum ecosystem. The main goal was to
relieve execution-layer nodes from the need to re-execute every new block by
means of a ZkVM. To achieve this, lattice-based cryptography was used, which,
unlike elliptic-curve-based approaches, is post-quantum secure and, unlike
hash-based approaches, can produce proofs in the order of tens of kilobytes
rather than hundreds or thousands.

The thesis succeeded in designing the architecture and partially implementing a
prototype lattice-based ZkVM. On the 100th Fibonacci benchmark, the prototype
produced a proof in 688.24 seconds using 12.34 GB of RAM. In a longer EVM
experiment, the system ran for approximately 9 hours and used 13.7 GB of RAM.
The results show that the approach is currently slower than established ZkVM
solutions, but they also suggest that this negative result is not caused by a
fundamental limitation and that further optimization is needed. In terms of
proof size, the post-quantum secure proofs are 4 to 5 times smaller than the
compared solutions.

\newpage{}\thispagestyle{empty}

\emptypage
